\documentclass[acmsmall,review=false]{acmart}

\settopmatter{printacmref=false}
\renewcommand\footnotetextcopyrightpermission[1]{}
\setcopyright{none}
\acmConference{}{}{}
\acmYear{2026}
\acmDOI{}
\acmISBN{}
\pagestyle{plain}

\usepackage{booktabs}
\usepackage{enumitem}
\usepackage{url}

\title{Detecting Deepfake Videos with Vision-Language Models: An Application of SmolVLM2}

\author{Likhitha Sindhu Geddam}
\affiliation{%
  \institution{Research Internship, IIIT Bangalore}
  \country{India}
}

\begin{document}

\begin{abstract}
Deepfake video detection systems are typically trained as opaque binary
classifiers, producing a real/fake verdict without any indication of \emph{why}
that verdict was reached. This lack of explainability limits their usefulness
in settings where a human still needs to verify or act on the model's output.
This report describes the fine-tuning of a compact vision-language model,
SmolVLM2-256M-Instruct~\citep{marafioti2025smolvlm}, into a multi-task deepfake
detector that jointly predicts a real/fake label and generates a natural-language
description of the visual artifacts supporting that prediction. The model is
adapted using Low-Rank Adaptation (LoRA)~\citep{hu2021lora} on both the attention
and feed-forward projections of the language backbone, alongside a lightweight
classification head and a contrastive vision-text alignment objective. Training
data is drawn from a compilation of established deepfake benchmarks, including
FaceForensics++~\citep{rossler2019faceforensics}, Celeb-DF~\citep{li2020celebdf},
and CREMA-D~\citep{cao2014cremad}. We report the engineering pipeline used to make
training feasible under a constrained compute and storage budget, how the
dataset's natural real/fake class imbalance is addressed in the classification
loss, and preliminary classification and generation results.
\end{abstract}

\maketitle

\section{Introduction}

The rapid advancement of generative AI tools has significantly lowered the
barrier to creating convincing deepfake videos, placing capabilities once
limited to specialized research groups into the hands of casual users. The
resulting circulation of such content online poses serious challenges,
including the spread of misinformation, reputational harm, and an erosion of
trust in visual media as reliable evidence~\citep{hondru2025exddv}. Most
deepfake videos fall into one of a few manipulation categories --- identity
swapping, expression reenactment, or fully AI-synthesized generation --- each
leaving behind subtle but distinct visual traces~\citep{rossler2019faceforensics}.

This project takes a different approach to detecting such videos. Rather than
treating detection as a purely binary classification problem, as the majority
of existing deepfake detectors do~\citep{li2020celebdf}, we fine-tune SmolVLM2,
a compact vision-language model, to both classify a video as real or fake
\emph{and} generate a concise natural-language explanation of the specific
visual artifacts supporting that decision --- describing inconsistencies in
texture, lighting, eye movement, facial motion, or boundary edges around the
face. This gives a human reviewer something actionable to verify, rather than
a bare label with no justification.

SmolVLM2 was chosen specifically for this task because it is designed to be
lightweight and computationally efficient: at only 256 million parameters, it
is built to run in resource-constrained settings where larger vision-language
models are impractical~\citep{marafioti2025smolvlm}. This made it a natural fit
for a project explicitly aimed at establishing whether a multi-task
detection-and-explanation model could be trained and run without requiring
large-scale computing infrastructure.

The remainder of this report is organized as follows. Section~\ref{sec:related}
reviews related work on deepfake detection datasets, explainability, and
parameter-efficient fine-tuning. Section~\ref{sec:dataset} describes the
training data, including a class-imbalance issue discovered and corrected
during the project. Section~\ref{sec:method} describes the model architecture
and multi-task training objective. Section~\ref{sec:infra} documents the
training infrastructure and engineering challenges encountered.
Section~\ref{sec:results} reports experimental results, and
Section~\ref{sec:conclusion} concludes.

\section{Related Work}
\label{sec:related}

\subsection{Deepfake Detection Datasets}

FaceForensics++~\citep{rossler2019faceforensics} is among the most widely used
benchmarks for facial forgery detection, comprising 1,000 real videos sourced
from YouTube, each manipulated using four representative techniques: DeepFakes,
Face2Face, FaceSwap, and NeuralTextures, spanning both identity-swap and
expression-reenactment manipulation families. Celeb-DF~\citep{li2020celebdf}
was introduced specifically to address a visual-quality gap in earlier
datasets: its authors observed that deepfakes circulating in the wild were
substantially higher quality than those in existing benchmarks, and released
5,639 refined DeepFake videos of celebrities to better reflect this. Both
datasets remain standard training and evaluation sources for forgery detection
research. CREMA-D~\citep{cao2014cremad}, while not itself a deepfake dataset,
is a widely used audio-visual corpus of emotionally expressive facial and vocal
recordings from 91 actors, and is used in this project as a source of diverse,
naturalistic real facial video content.

\subsection{Explainable Deepfake Detection}

\citet{hondru2025exddv} introduce ExDDV, the first dataset explicitly designed
for explainable deepfake video detection, pairing roughly 5.4K real and fake
videos with human-written textual artifact descriptions and spatial click
annotations localizing the described artifacts. Their experiments, which
fine-tune and in-context-prompt several vision-language models, find that
combining textual and spatial supervision is necessary to obtain models capable
of both localizing and describing observed artifacts, and that description
quality alone (without spatial grounding) remains a meaningful challenge for
current models. This motivates our decision to frame artifact explanation as a
learned generative task rather than a post-hoc justification step, while
acknowledging that our setup, lacking spatial click supervision, targets a
strictly text-only variant of this explainability goal.

\subsection{Parameter-Efficient Fine-Tuning}

Adapting large pretrained models to a downstream task traditionally requires
updating all model parameters, which is memory- and compute-intensive and
produces a full-size checkpoint per task. \citet{hu2021lora} propose Low-Rank
Adaptation (LoRA), which freezes the pretrained weight matrices and instead
learns a pair of low-rank matrices whose product is added to each targeted
weight matrix at inference time. Because the rank of these matrices is chosen
to be much smaller than the original matrix dimensions, the number of trainable
parameters is reduced by orders of magnitude relative to full fine-tuning, while
empirically retaining competitive downstream performance. This project applies
LoRA to both the attention projections (query, key, value, output) and the
feed-forward projections (gate, up, down) of the SmolVLM2 language backbone,
while keeping the vision encoder frozen and the multimodal connector and
task-specific heads fully trainable.

\subsection{Compact Vision-Language Models}

\citet{marafioti2025smolvlm} introduce the SmolVLM family, a set of
vision-language models explicitly engineered for resource-efficient inference,
released at 256M, 500M, and 2.2B parameter scales. The authors report that
architectural choices around visual tokenization and data curation allow even
their smallest 256M-parameter variant to outperform substantially larger prior
models on a memory-normalized basis, making it a practical backbone for
on-device or low-resource fine-tuning scenarios such as this one. SmolVLM2
extends this family with native video understanding, processing a sequence of
sampled frames as an interleaved image sequence, which is the mechanism this
project relies on for temporal (multi-frame) input.

\section{Dataset}
\label{sec:dataset}

\subsection{Composition}

The training corpus is compiled from multiple established sources described
above --- FaceForensics++, Celeb-DF, and CREMA-D --- combined with additional
AI-generated fake content, and hosted as a private dataset repository on the
Hugging Face Hub.\footnote{\url{https://huggingface.co/datasets/LikhithaGeddam/fakevideodetection-private}}
Each video is labeled \texttt{real} or \texttt{fake} and organized into
\texttt{train}, \texttt{val}, and \texttt{test} splits. Table~\ref{tab:splits}
gives the per-split class counts.

\begin{table}[h]
\centering
\caption{Dataset split composition by class.}
\label{tab:splits}
\begin{tabular}{lccc}
\toprule
\textbf{Split} & \textbf{Real} & \textbf{Fake} & \textbf{Total} \\
\midrule
Train & 6,532 & 9,634 & 16,166 \\
Val   & 1,400 & 1,904 & 3,304 \\
Test  & 1,400 & 2,065 & 3,465 \\
\bottomrule
\end{tabular}
\end{table}

\subsection{Caption Generation and Class Weighting}

To supervise the natural-language explanation task, each training video
requires an associated artifact description. Captions were generated by
running a pretrained deepfake detection model
(\texttt{prithivMLmods/Deep-Fake-Detector-v2-Model}, a Vision Transformer-based
classifier) over sampled frames of each video, recording its prediction,
per-frame confidence, and an attention-rollout-derived description of the
facial region most responsible for the prediction. This process was run to
completion across the full training pool, giving all 6,532 real and all 9,634
fake training videos (Table~\ref{tab:splits}) a detector-derived caption, with
uniform coverage across both classes.

Although caption coverage is complete and uniform, the underlying video corpus
is itself class-imbalanced: the training split is 40.4\% real and 59.6\% fake
by video count (Table~\ref{tab:splits}). Left unaddressed, this imbalance
could let the classification head partially rely on the class prior rather
than learning a genuine decision boundary. To correct for this, the
classification loss uses a dataset-level fixed class weight computed once
from the true training-set class counts, rather than a per-batch reweighting
scheme whose effective weighting can vary with the specific class composition
of each mini-batch. This fixed weighting is applied identically at every
training step regardless of batch composition (Section~\ref{sec:method}).

\section{Multimodal Signal Design}
\label{sec:signal}

Before describing the model architecture in detail (Section~\ref{sec:method}),
we summarize the multimodal signal the system is built around: what features
it detects on, how inputs are constructed, and how the two modalities are
represented and fused.

\paragraph{Features used for detection.} The model does not use any hand-designed
artifact detector. Instead, visual artifacts (compositing boundaries,
texture/lighting inconsistencies, unnatural motion) are learned implicitly by
the vision encoder and connector, supervised jointly by (i) a binary real/fake
classification signal and (ii) natural-language artifact descriptions that
name the specific inconsistency observed (Section~\ref{sec:dataset}). No audio,
transcript, or external metadata signal is used --- the system is
visual-and-caption-only.

\paragraph{Input construction.} Each training example consists of a fixed
number of frames sampled per video clip (stratified temporal sampling with
random jitter, Section~\ref{sec:method}), paired with a text prompt requesting
a real/fake judgment and explanation, and a target text sequence containing
the ground-truth label and artifact description. No face-cropping or alignment
preprocessing is applied; frames are passed to the vision encoder at their
sampled resolution and the model must localize any facial region itself.

\paragraph{Input representation and fusion.} Frames are tokenized by the
SigLIP-based vision encoder and projected into the language model's embedding
space by a pixel-shuffle connector, then interleaved with text tokens as a
single sequence --- an image-text pair representation rather than a separate
late-fusion classifier. Classification and contrastive alignment
(Section~\ref{sec:method}) both operate on the same connector output captured
via a forward hook during this single forward pass, rather than a second,
detached vision pass; this makes the fusion embedding-level rather than
decision-level (late) fusion, since the classification decision is read
directly off the fused visual representation rather than combined post-hoc
with a separate text-based decision.

\section{Method}
\label{sec:method}

\subsection{Backbone Architecture}

The base model is SmolVLM2-256M-Instruct~\citep{marafioti2025smolvlm}, comprising
a SigLIP-based vision encoder, a pixel-shuffle connector projecting visual
features into the language model's embedding space, and a SmolLM2-derived causal
language decoder. For video input, a fixed number of frames are sampled per clip
using stratified temporal sampling with a small random jitter, encouraging
coverage across the full clip duration while retaining some sampling diversity
across epochs.

\subsection{Multi-Task Heads}

On top of the shared backbone, we attach two lightweight task-specific heads:

\begin{itemize}[nosep]
    \item \textbf{Classification head}: a two-layer MLP (576 $\rightarrow$ 512
    $\rightarrow$ 2, with ReLU and dropout) applied to the mean-pooled visual
    connector output, producing a real/fake logit pair.
    \item \textbf{Contrastive alignment head}: linear projections of the pooled
    visual features and the final-token text hidden state into a shared
    256-dimensional embedding space, trained with a symmetric InfoNCE-style
    contrastive loss (temperature $= 0.07$) to encourage visual and textual
    representations of the same video to align.
\end{itemize}

A key implementation detail is that both auxiliary heads consume visual
features captured via a forward hook on the connector module \emph{during} the
same forward pass used for language modeling, rather than via a second,
gradient-detached vision encoder pass. This ensures classification and
contrastive gradients propagate back through the connector (and, when unfrozen,
the vision encoder), rather than being computed on frozen, non-differentiable
features --- a design choice made after observing that a detached-feature
variant caused the auxiliary losses to plateau early without meaningfully
shaping the visual representation.

\subsection{Parameter-Efficient Adaptation}

Rather than fine-tuning all backbone parameters, we apply LoRA
\citep{hu2021lora} (rank 32, $\alpha=64$) to the query, key, value, and output
attention projections as well as the gate, up, and down feed-forward
projections of the language decoder. The vision encoder is frozen; the
multimodal connector and both task heads remain fully trainable. This
configuration was chosen to balance adaptation capacity against the
overfitting risk and compute cost of full fine-tuning on a moderately sized
video corpus.

\subsection{Training Objective}

The total training loss is a weighted combination of three terms:
\begin{equation}
\mathcal{L} = \lambda_{\text{cls}} \mathcal{L}_{\text{cls}} + \lambda_{\text{lm}} \mathcal{L}_{\text{lm}} + \lambda_{\text{contr}}(t) \, \mathcal{L}_{\text{contr}}
\end{equation}
where $\mathcal{L}_{\text{cls}}$ is class-weighted cross-entropy on the
real/fake prediction (weights computed once from full training-set class
counts, per Section~\ref{sec:dataset}), $\mathcal{L}_{\text{lm}}$ is the
standard causal language-modeling loss on the generated explanation text (with
the prompt/header tokens masked out of the loss via direct token-ID
matching of the assistant turn boundary), and $\mathcal{L}_{\text{contr}}$ is
the contrastive alignment loss described above, whose weight $\lambda_{\text{contr}}(t)$
is ramped linearly from zero over the first portion of training rather than
applied at full strength from step one, since a very small effective batch
size gives an unreliable pool of contrastive negatives early in training. A
label-smoothing term ($\epsilon = 0.05$) is applied to the classification loss.

\section{Training Infrastructure}
\label{sec:infra}

Beyond the modeling choices above, a substantial portion of this project
involved making training \emph{feasible at all} under real-world resource
constraints, which we document here as it materially affected the scope and
schedule of experiments.

\subsection{Storage-Constrained Data Access}

The full video corpus, at native resolution, exceeds the available local
storage in a free-tier cloud notebook environment (Google Colab) by roughly an
order of magnitude, so storing or repeatedly re-downloading raw video was not
viable. We addressed this with a two-stage, precompute-then-train pipeline
rather than reading video directly during training:

\begin{enumerate}[nosep]
    \item \textbf{One-time frame precomputation}: a separate script
    (\texttt{precompute\_frames\_streaming.py}) streams each video once from
    the Hugging Face Hub-hosted dataset repository (Section~\ref{sec:dataset}),
    extracts a fixed pool of 8 frames per video, and writes them to local disk
    as JPEGs under a \texttt{split/\{real,fake\}/<video\_stem>/} directory
    structure, together with \texttt{train}/\texttt{val}/\texttt{test} JSON
    manifests that reference frame paths rather than a video path. The source
    video is not retained after its frames are extracted.
    \item \textbf{Frames-only training}: the training script reads exclusively
    from this precomputed frame pool and never touches a video file or a video
    decoding library. \texttt{VideoFakeClueDatasetFrames} loads the relevant
    JPEGs directly with PIL, so the \texttt{decord}-based video decoding used
    in earlier iterations of the pipeline is removed entirely from the
    training path.
\end{enumerate}

Because only JPEG frames are stored rather than full videos, this keeps the
persistent disk footprint under 10\,GB, well within a free-tier notebook
environment's budget, without needing per-video Hub streaming during training
itself. Per-epoch sampling variety is preserved despite frames being extracted
only once: each \texttt{\_\_getitem\_\_} call still selects a jittered subset
of \texttt{num\_frames} (4, by default) from the saved 8-frame pool and applies
random flip/brightness augmentation, so the same video contributes a
different, augmented frame subset on different epochs even though the
underlying JPEGs on disk are fixed. Before training begins, a resolution check
verifies that every entry's referenced frame files are actually present on
disk, reporting any gap as an explicit count and percentage rather than
silently degrading data quality with placeholder frames.

\subsection{Checkpointing and Fault Tolerance}

Because training runs extend over many hours on infrastructure subject to
session timeouts, disconnects, and preemption, checkpointing was treated as a
first-class concern rather than an afterthought. Model, optimizer, and
scheduler state are checkpointed to persistent storage at the end of every
epoch, and training can resume from an arbitrary checkpoint rather than
restarting. A separate, continuously-updated ``best'' checkpoint is maintained
based on held-out classification accuracy (Section~\ref{sec:results}),
independent of the most recent epoch checkpoint, so that a partially-completed
or subsequently-degraded epoch does not overwrite the best model observed so
far.

\subsection{Reproducibility Under a Changing Environment}

A secondary, less anticipated challenge was environment drift: the training
script was, at different points, affected by a background library update
(\texttt{transformers}) that removed a previously-standard training argument
in a major version release. Rather than pin a single environment version
indefinitely, the training script was made defensive to this class of failure:
it inspects the installed library's accepted arguments at runtime and
translates or drops unsupported configuration values with an explicit warning,
rather than failing on an unrecognized keyword argument mid-run.

\section{Experimental Setup}
\label{sec:setup}

\begin{table}[h]
\centering
\caption{Key training hyperparameters.}
\label{tab:hparams}
\begin{tabular}{ll}
\toprule
\textbf{Hyperparameter} & \textbf{Value} \\
\midrule
Base model & SmolVLM2-256M-Instruct \\
LoRA rank / $\alpha$ & 32 / 64 \\
LoRA target modules & q,k,v,o\_proj, gate,up,down\_proj \\
Frames per clip & 4--6 (stratified + jitter) \\
Per-device batch size & 4 \\
Gradient accumulation steps & 2--4 \\
Effective batch size & 8--16 \\
Optimizer & AdamW ($\beta_1{=}0.9$, $\beta_2{=}0.999$) \\
Peak learning rate & $5\times10^{-5}$ \\
LR schedule & Cosine, with warmup \\
Precision & bfloat16 \\
Epochs & 3 \\
\bottomrule
\end{tabular}
\end{table}

Head parameters (classification and contrastive heads) are trained at a higher
relative learning rate than the base LoRA parameters, reflecting that these
components are randomly initialized and have no pretrained prior to preserve.
The dataset's pre-defined \texttt{val} split (3,304 videos, disjoint from the
training videos actually seen by the optimizer; Table~\ref{tab:splits}) is used
to track real/fake classification accuracy at the end of every epoch,
independent of the training-batch accuracy logged during optimization, which
we found could be misleadingly high on individual batches without reflecting
genuine generalization.

\section{Results}
\label{sec:results}

\subsection{Classification Performance}

Classification metrics below were obtained by reloading each epoch's saved
checkpoint (\texttt{checkpoint-910}, \texttt{checkpoint-1820}, and
\texttt{checkpoint-2730}, corresponding to epochs 1--3) and evaluating on a
stratified subsample of 400 videos drawn from the validation split (166 real,
234 fake), preserving the split's original class ratio. This post-hoc
evaluation was necessary because the training-time validation callback only
logged accuracy to console output rather than persisting it, and did not
compute per-class precision/recall at all; the callback has since been
extended to log all three metrics to disk for future runs.

\begin{table}[h]
\centering
\caption{Validation classification performance by epoch (stratified 400-video subsample of the validation split).}
\label{tab:results}
\begin{tabular}{lccc}
\toprule
\textbf{Epoch} & \textbf{Accuracy} & \textbf{Precision (real)} & \textbf{Recall (real)} \\
\midrule
1 & 91.8\% & 99.3\% & 80.7\% \\
2 & 91.5\% & 97.1\% & 81.9\% \\
3 & 91.8\% & 98.5\% & 81.3\% \\
\bottomrule
\end{tabular}
\end{table}

Given the class imbalance documented in Section~\ref{sec:dataset}, we report
precision and recall for the minority (real) class specifically, rather than
overall accuracy alone, since overall accuracy can be misleadingly high on an
imbalanced split even for a classifier that has not learned a meaningful
decision boundary. Precision on the real class stays consistently above 97\%
across all three epochs, indicating the model rarely misclassifies a fake
video as real; recall plateaus around 81\%, indicating roughly one in five
real videos is still misclassified as fake. Accuracy, precision, and recall
are all essentially flat across epochs 1--3, suggesting the classification
head converges early and additional epochs primarily benefit the
explanation-generation objective rather than the classification objective.

\subsection{Ablation Studies}
\label{sec:ablations}

We ran three ablations against a matched baseline, each isolating one design
choice: (i) \emph{single\_frame}, sampling 1 frame per clip instead of 4, to
test the contribution of temporal information; (ii) \emph{no\_contrastive},
setting the contrastive vision-text alignment weight to 0 instead of its
ramped schedule (Section~\ref{sec:method}), to test the contribution of the
alignment objective; and (iii) \emph{lower\_lr}, a peak learning rate of
$1\times10^{-5}$ instead of $5\times10^{-5}$. Because a full 3-epoch run per
variant is not feasible on a single Colab GPU, all four runs (including the
baseline) use a fixed, short 350-step budget with a freshly initialized model,
and are evaluated on the same stratified 200-video subsample of the validation
split (identical across variants, disjoint from the subsamples used in
Table~\ref{tab:results}).

\begin{table}[h]
\centering
\caption{Ablation results (350-step budget, evaluated on a shared 200-video validation subsample).}
\label{tab:ablations}
\begin{tabular}{lcccc}
\toprule
\textbf{Variant} & \textbf{Accuracy} & \textbf{Precision (real)} & \textbf{Recall (real)} & \textbf{Final train loss} \\
\midrule
baseline       & 91.0\% & 100.0\% & 78.3\% & 12.09 \\
single\_frame  & 91.0\% & 100.0\% & 78.3\% & 11.83 \\
no\_contrastive& 91.0\% & 100.0\% & 78.3\% & 11.91 \\
lower\_lr      & 89.0\% &  94.2\% & 78.3\% & 23.37 \\
\bottomrule
\end{tabular}
\end{table}

The learning-rate ablation shows a clear, expected effect: at a 5x lower peak
LR, the fixed 350-step budget is not enough to converge, leaving final
training loss roughly double that of the other three variants and a
corresponding drop in accuracy and real-class precision. This confirms the
step budget itself is a meaningful constraint on these ablation results,
not just a convenience.

The \emph{single\_frame} and \emph{no\_contrastive} variants, by contrast,
converged to \emph{exactly} the same classification accuracy, precision, and
recall as the baseline, despite visibly different per-step loss values
throughout training and, for
\emph{single\_frame}, roughly 4.6x faster wall-clock training (5.5 s/step vs.
25.1 s/step, since a single frame is far cheaper to encode than four). We do
not read this as evidence that frame count and contrastive alignment have no
effect on the model; rather, we consider it likely that (a) at only 350 steps
and 200 evaluation videos, classification accuracy is not a sensitive enough
metric to detect a difference this small, and (b) the classification decision
may be dominated by per-frame, low-level artifacts (compression, resolution,
lighting) that a single frame already captures, while temporal information
and cross-modal alignment may matter more for the quality of the generated
\emph{explanation} than for the binary classification decision itself. We
flag this as a specific, motivated direction for future ablation work: repeat
this comparison using generation-quality metrics (SemSim, ROUGE-L, BLEU-1) in
addition to classification accuracy, and with a longer step budget and larger
eval set to increase sensitivity. One immediate practical takeaway does hold
regardless: single-frame training reaches the same classification accuracy at
roughly 4.6x lower wall-clock cost, which is useful for fast iteration during
hyperparameter search even if the final model should retain multiple frames
for explanation quality.

\subsection{Qualitative Explanation Generation}

Representative generated explanations, sampled periodically during training
against held-out examples, are shown below.

\begin{quote}
\textit{Ground truth}: ``This video is FAKE. Micro-expressions are absent; the
face shows only coarse motion without fine muscle detail. The face appears to
lag behind the head motion by one or two frames.''\\
\textit{Generated}: ``This video is FAKE. The face has an unnaturally uniform
skin tone without the natural micro-variation of real skin. Lighting direction
on the face is inconsistent with the environmental light source.''
\end{quote}

The classification label and the specific artifact described agree with the
ground truth here, though the generated explanation identifies a different
artifact (skin tone / lighting) than the one a human annotator focused on
(motion lag), which is expected given each is one plausible artifact among
several present in a genuinely manipulated video.

A small held-out qualitative test (4 videos, 2 real and 2 fake, evaluated with
the fully fine-tuned model) gave a classification accuracy of 3/4 (75.0\%),
average semantic similarity of 0.838, average ROUGE-L of 0.260, and average
BLEU-1 of 0.392 between generated and ground-truth explanations. The one
misclassification and two label-disagreement cases from this same test are
examined in Section~\ref{sec:failure}.

\section{Failure Analysis}
\label{sec:failure}

We examine the three non-trivial cases from the 4-video held-out qualitative
test (Section~\ref{sec:results}) in more detail, since a 75\% accuracy result
on so few videos is more informative as individual case studies than as a
single aggregate number.

\paragraph{Misclassified real video.} On \texttt{train\_real\_5983.mp4} (ground
truth REAL), the classification head predicted FAKE with 82.0\% confidence ---
the one classification error in this test set. The generated explanation for
this video also contained a malformed label token: it began ``This video is
FEARFUL!'' rather than a well-formed REAL/FAKE judgment. We do not have enough
occurrences of this token corruption to characterize it statistically, but its
co-occurrence with the one classification error in this small test is worth
noting: it may indicate the model was in a genuinely uncertain state on this
example (consistent with the classification head's own 82.0\% rather than
higher confidence), where both heads produced lower-quality output together,
rather than the two failures being independent.

\paragraph{Classification/generation label disagreement.} In 2 of the 4 test
videos, the classification head's predicted label and the label stated in the
first sentence of the \emph{generated explanation text} disagreed with each
other, even though the classification head's prediction was itself correct in
both cases. On \texttt{train\_1281.mp4} (ground truth FAKE), the classification
head correctly predicted FAKE with 92.6\% confidence, but the generated
explanation began ``This video is Real.'' A symmetric case occurred on
\texttt{train\_real\_4670.flv} (ground truth REAL), where the classification
head correctly predicted REAL with 70.7\% confidence, but the generated
explanation began ``This video is FAKE.'' Unlike the misclassification above,
these are cases where the \emph{classification head} got the right answer but
the \emph{language-modeling head}'s stated label did not agree with it. This
indicates that, at least on this small sample, the two heads can disagree on
the surface label they each produce even when one of the two is correct --- a
direct consequence of the two heads being trained with separate loss terms
(Section~\ref{sec:method}) rather than any mechanism that enforces label
agreement between them.

\paragraph{Implication.} Together, these three cases point to the same
underlying gap: nothing in the current training objective ties the
classification head's decision to the label asserted by the generation head.
A consistency loss (e.g.\ penalizing disagreement between the two heads'
labels during training) or constrained decoding (forcing the generated label
token to match the classification head's argmax at inference time) are both
direct, implementable responses to this finding, and we list both as concrete
future work in Section~\ref{sec:conclusion}.

\section{Limitations}
\label{sec:limitations}

Several limitations should be noted. First, training used a reduced per-clip
frame count relative to what full-video temporal analysis would ideally use,
a direct consequence of memory and throughput constraints on the available
hardware. Second, the model lacks the spatial click-localization supervision
used by \citet{hondru2025exddv}, meaning it can describe but not point to a
specific image region; extending the architecture with a grounding head is a
natural direction for future work. Third, the classification head and the
generation head can disagree on the label each produces even when one is
correct, and a small number of generated explanations contained malformed
label tokens; see Section~\ref{sec:failure} for a detailed analysis of both
issues. Finally, results reported in Section~\ref{sec:results} reflect
training on a compiled multi-source dataset rather than a single canonical
benchmark, which should be considered when comparing against results reported
on FaceForensics++ or Celeb-DF alone.

\section{Conclusion and Future Work}
\label{sec:conclusion}

This project demonstrates that a compact, 256M-parameter vision-language model
can be adapted via parameter-efficient fine-tuning into a functioning
multi-task deepfake detector, jointly producing a classification decision and
a natural-language justification, under meaningfully constrained compute and
storage conditions. Beyond the modeling contribution, the project surfaced an
engineering finding we consider independently useful: a checkpoint- and
cache-persistent training pipeline substantially changes the practical
feasibility of long training runs on interruptible infrastructure. Future
work includes extending
supervision with spatial artifact localization following
\citet{hondru2025exddv}, adding a consistency loss or constrained decoding to
close the classification/generation label-disagreement gap identified in
Section~\ref{sec:failure}, evaluating cross-dataset generalization (training on
this compiled corpus and testing on a held-out canonical benchmark such as
FaceForensics++ alone), and scaling to the 500M or 2.2B SmolVLM2 variants once
additional compute is available.


\begin{thebibliography}{9}

\bibitem[Cao et al.(2014)]{cao2014cremad}
Cao, H., Cooper, D. G., Keutmann, M. K., Gur, R. C., Nenkova, A., \& Verma, R. (2014).
CREMA-D: Crowd-sourced emotional multimodal actors dataset.
\textit{IEEE Transactions on Affective Computing}, 5(4), 377--390.

\bibitem[Hondru et al.(2025)]{hondru2025exddv}
Hondru, V., Hogea, E., Onchis, D., \& Ionescu, R. T. (2025).
ExDDV: A new dataset for explainable deepfake detection in video.
\textit{arXiv preprint arXiv:2503.14421}. (Accepted at WACV 2026.)

\bibitem[Hu et al.(2021)]{hu2021lora}
Hu, E. J., Shen, Y., Wallis, P., Allen-Zhu, Z., Li, Y., Wang, S., Wang, L., \& Chen, W. (2021).
LoRA: Low-rank adaptation of large language models.
\textit{arXiv preprint arXiv:2106.09685}.

\bibitem[Li et al.(2020)]{li2020celebdf}
Li, Y., Yang, X., Sun, P., Qi, H., \& Lyu, S. (2020).
Celeb-DF: A large-scale challenging dataset for deepfake forensics.
In \textit{Proceedings of the IEEE/CVF Conference on Computer Vision and Pattern Recognition (CVPR)}, 3207--3216.

\bibitem[Marafioti et al.(2025)]{marafioti2025smolvlm}
Marafioti, A., Zohar, O., Farr\'e, M., Noyan, M., et al. (2025).
SmolVLM: Redefining small and efficient multimodal models.
\textit{arXiv preprint arXiv:2504.05299}.

\bibitem[R\"ossler et al.(2019)]{rossler2019faceforensics}
R\"ossler, A., Cozzolino, D., Verdoliva, L., Riess, C., Thies, J., \& Nie{\ss}ner, M. (2019).
FaceForensics++: Learning to detect manipulated facial images.
In \textit{Proceedings of the IEEE/CVF International Conference on Computer Vision (ICCV)}, 1--11.

\bibitem[Unleashing(2026)]{unleashing2026vlsemantics}
Unleashing vision-language semantics for deepfake video detection.
\textit{arXiv preprint arXiv:2603.24454}, 2026.

\end{thebibliography}

\end{document}
